\documentclass[12pt]{article}
\usepackage[utf8]{inputenc}
\usepackage[T1]{fontenc}
\usepackage{lmodern}
\usepackage{microtype}
\usepackage[a4paper, margin=3cm]{geometry}
\usepackage{parskip}

\date{\today}

\begin{document}

\raggedright

{\small\the\year}

\bigskip

Dear Sahaj family,

\bigskip

I'm writing you to express my intention to make a change concerning our communication tools in the collective.
We live in times where connectivity is as easy as ever and where everything is at our finger tips.
Any information that is required --- it can be googled, "chatgpt'd", or I can call or text anyone in the world to ask for it, immediately.
The power of it does not need any explanation.
However, what does it do to our attention or the quality of our communication?
If there is such a low barrier to reach out to someone, then what's the need to think twice what I'm going to ask?
Just like the beautiful video that I just saw and that I can just easily forward to a few hundred people without having to give them any context why I care about sharing it to them.
And if I'm available all the time for everyone, then when do I focus my attention without interruption?
Also, how do I filter all the content that I receive on a daily basis that is simply too much to take in all?
Now, these are all problems that we can relate to and mostly found our fix already.
Unfortunately though, these are not even the gravest problems that this connectivity industry causes.
As we all know there is a reason why these services, also referred to as \textit{Social Media}, are free --- because we are the product.
Or more precisely, the data we generate, is.
I know what you're thinking now: "Ahh okay, now I get what this letter is about. He's also one of those data privacy freaks."
That might be, but this is not what this message is about today.
Hear me out.
So, we all know that Google, Facebook \& Co. collect data through their service.
Things like, which device you're on, which websites you visit, how long you visit them, then who of your friends you message when and how often.
All this data is anonymized.
That means, they're not interested in the exact names.
What's important is how they can analyze your internet behaviour and your social network.
And that's particularly interesting for marketing companies.
So, Google \& Co. collect a bunch of data --- most of it not very interesting in isolation --- and sells it to marketing companies that then help their clients market their products more effectively.
E.g., Philips brings out new airfryer and wants to show their new amazing product to all interested customers.
They launch an online ad and use internet traffic data to find out, who looked at airfryers recently, or who even bought one, or who fits the profile of someone who buys airfryers.
Maybe a middle aged woman, from a middle class background and above average socio-economic status.
All this information is easily accessible through our internet traffic footprint.
Depending on our surfing behaviour, one can find out all these things.
Additionally, what works wonders in marketing is the social effect.
I.e., when my friend buys something I'm much more likely to buy it as well.
And this is where the social network information comes into play.
If an online ad provider knows that you bought an airfryer recently, it will target all your friends for the next weeks and show them a lot of airfryer ads.
Eventually, they're primed and start looking for airfryers.
This is when people often say: "That's crazy, last week I was talking about airfryers with my friend, and now I'm getting all these ads. It's almost like my phone is listening to us!"
That sound scary, but is exaggerated.
The joke is that they don't even need to listen to your conversation with your friend to know that you talked about airfryers.
It's enough to know that one of you bought one recently and that you are good friends because you send each other messages on a regular basis and talk to each other frequently on the phone.
The amount of reels you send each other on instagram perhaps is the gauge for how close you are.
Okay, so what?
Great, those companies collect our data and sell more online ads to companies.
What's the need to make a big fuss about it?
Our economy is dependent on consumption and it's kind of useful that the things I need to buy are easily available for me to find.
The reason I'm writing this message is not because of that reason, but because of another thing that those companies can do --- manipulating people's political choice.
There is proof that companies like "Cambridge Analytica" helped political campaigns, like that of Brexit and of Donald Trump to win through manipulating voters by targeting them specifically and also by using fake news to change their sentiment.
The implications of this are huge.
Democracy as such is not able to be held up in such a world, and I believe it is essential that the power of the internet industry is reduced.
For that to happen we can wait for governments to make up new laws to protect the data of users, or we can start acting already because we actually have a lot of power ourselves.
If you don't believe that then ask yourself who is making Google, Facebook, and Amazon rich.
It's you!
You're the one using services of Google, Facebook, Instagram, Whatsapp, and buying things on Amazon.
And that in turn means we can also play a part making them less and rich and powerful by using their services less.
My point is not to abandon the internet entirely, or to strictly condemn any of those companies.
Instead we can be more conscious of the services we use and replace some of them with other options.
And today I therefore want to propose we move our Sahaja Yoga communication away from Whatsapp.
I believe Sahaja Yogis should be an example for society and so I think we should be pioneers when it comes to doing the right.
Apart from that Whatsapp is a terrible messenger for teams and organizations.
It was never meant to fulfill this role, but was to function as a direct messenger for small and quick communication between individuals and small groups.
However, because of convenience we want everything to be in Whatsapp --- which is very understandable.
Nevertheless, I see so many shortcomings in our communication channels.
It is difficult to send important information to everyone without it getting lost in other messages.
Also, the ease of writing something leads to an abundance of out-of-context forwards in big chats.
Another point is the exploding number of groups that are difficult to keep track of, paired with the mixing of all kinds of chats in the inbox.
If we are serious about changing the world for the better we should consider some alternatives, like \textit{Slack}.
Not only to stop supporting the data industry, but also to enable a more effective collaboration among us Yogis --- because our efficacy matters to the world!

\bigskip

Sincerely,\\
Ganesh

\end{document}
